JWST has identified exceptionally compact, luminous galaxies at redshifts
$z > 7$. To understand the nature of these sources, we must analyze the
observational data and test whether they are consistent with basic
astrophysical constraints. The table below lists the constants and parameters
used in this problem.

\begin{center}
\begin{tabular}{p{0.55\linewidth}c}
\toprule
Star formation rate (SFR) to Luminosity conversion factor &
  $k = 1.0 \times 10^{10}\,\mathrm{L_\odot\,M_\odot^{-1}\,yr}$ \\
Effective opacity of free electrons &
  $\kappa_{\mathrm{e}} = 0.04\,\mathrm{m^2\,kg^{-1}}$ \\
Effective opacity of the dusty interstellar medium &
  $\kappa_{\mathrm{dust}} = 1.5\,\mathrm{m^2\,kg^{-1}}$ \\
Optical-to-bolometric flux correction factor &
  $BC = F_{bol}/F_{opt} = 9.26$ \\
\bottomrule
\end{tabular}
\end{center}

\medskip\noindent\textbf{Part 1 (13 points): Cosmological Surface Brightness Dimming}

Let $L$ be the intrinsic bolometric luminosity and $R$ be the intrinsic
physical radius of a galaxy. The intrinsic surface luminosity density is
defined as
\[ \Sigma_{\mathrm{int}} = \frac{L}{\pi R^2} \]
A telescope observes this galaxy subtending with a solid angle
$\Omega_{\mathrm{obs}}$ and measures a flux $F_{\mathrm{obs}}$. The observed
surface brightness is
\[ \Sigma_{\mathrm{obs}} = \frac{F_{\mathrm{obs}}}{\Omega_{\mathrm{obs}}} \]
In cosmology, the luminosity distance $(D_L)$ is defined via
\[ F_{\mathrm{obs}} = \frac{L}{4\pi D_L^2} \]
The angular diameter distance $(D_A)$ is defined so that the observed angular
radius is
\[ \theta = \frac{R}{D_A} \]

\begin{description}
\item[(a)] (10 points) Derive the relation between $D_L$, $D_A$ and redshift
  $z$.
\item[(b)] (3 points) Derive an expression for $\Sigma_{\mathrm{obs}}$ as a
  function of $\Sigma_{\mathrm{int}}$ and redshift $z$.
\end{description}

\medskip\noindent\textbf{Part 2 (8 points): Dissecting JWST Data}

JWST observed an ultraluminous galaxy at a redshift of $z = 7.26$.
\begin{itemize}
\item The observed rest-frame optical flux at Earth (attenuated by cosmic
  dust) is $F_{\mathrm{obs,opt}} = 5.72 \times 10^{-19}\,\mathrm{W\,m^{-2}}$.
\item The dust extinction derived from the spectrum is
  $A_V = 2.79\,\mathrm{mag}$.
\item The measured angular radius of the source on the image is
  $\theta = 0.05\,\mathrm{arcsec}$.
\end{itemize}
Calculate the intrinsic bolometric surface luminosity density of this galaxy,
$\Sigma_{\mathrm{int}}$, in units of $\mathrm{L_\odot\,kpc^{-2}}$.

\medskip\noindent\textbf{Part 3 (18 points): Eddington Dynamical Model for the Star-Forming Disk}

Assume the entire intrinsic radiative energy of this galaxy originates from
extreme star formation in an optically thick disk of cold gas and dust.
Radiation from young stars generates tremendous pressure that drives cosmic
dust particles outward, perpendicular to the disk midplane. Due to the extreme
density, gas and dust are perfectly dynamically coupled via collisions. The
radiation pressure acting on the dust effectively drags the entire gas mass
along as a single unified fluid. Simultaneously, gravity pulls this fluid back
toward the midplane. Assume the gravitational field of the disk is governed by
the total dynamical surface mass density, given as
$\Sigma_{\mathrm{tot}} = 2.3 \times 10^{10}\,\mathrm{M_\odot\,kpc^{-2}}$.

\begin{description}
\item[(a)] (6 points) $\Sigma_{\mathrm{int\_max}}$ denotes the disk's
  Eddington surface luminosity density. Prove the acceleration driven by
  radiation pressure is
  $g_{\mathrm{rad}} = \dfrac{\kappa_{\mathrm{dust}}\Sigma_{\mathrm{int\_max}}}{c}$
  and the acceleration governed by the disk's gravity is
  $g_z = 2\pi G \Sigma_{\mathrm{tot}}$.
\item[(b)] (7 points) Calculate the SFR surface density at Eddington limit
  $(\Sigma_{\mathrm{Edd}})$ of this disk (in units of
  $\mathrm{M_\odot\,yr^{-1}\,kpc^{-2}}$) such that it remains gravitationally
  bound against a radiation-driven blowout.
\item[(c)] (5 points) Calculate the ratio
  $\Sigma_{\mathrm{SFR\_actual}}/\Sigma_{\mathrm{Edd}}$, where
  $\Sigma_{\mathrm{SFR\_actual}}$ is the actual SFR surface density of this
  galaxy. If the observed luminosity is purely stellar in origin, is this
  system dynamically stable?
\end{description}

\medskip\noindent\textbf{Part 4 (11 points): Clues from High-Resolution Spectroscopy}

Astronomers analyzed the NIRSpec spectrum and discovered that the emission
predominantly originates from a hot, highly ionized plasma, where continuous
recombination processes emit Hydrogen lines exhibiting anomalous broadening
with a velocity dispersion of $v = 3400\,\mathrm{km\,s^{-1}}$. Furthermore,
the radius of this emission region is indirectly estimated to be merely
$R_{\mathrm{em}} = 0.18\,\mathrm{pc}$.

\begin{description}
\item[(a)] (4 points) Calculate the dynamical mass $M$ confined within this
  region (in $\mathrm{M}_\odot$). Based on your result, choose the specific
  type of astrophysical object.
  \begin{description}
  \item[A.] A dwarf galaxy
  \item[B.] A massive star
  \item[C.] A super massive black hole
  \item[D.] A globular cluster
  \end{description}
\item[(b)] (7 points) Quantitatively evaluate the radiative stability of the
  system with respect to its observed luminosity. Is this system radiatively
  stable?
\end{description}
